\section{Parametrische Darstellung statt Abtastung}
\subsection{Motivation}
Klassische Diskretisierung stellt Krümmung und Überhöhung durch Stichproben
\(\kappa_i\approx\kappa(s_i)\), \(\phi_i\approx\phi(s_i)\) dar. Daraus entstehen
hochdimensionale Probleme mit schwacher Struktur und geringer
Interpretierbarkeit. AIM verwendet stattdessen parametrisierte Funktionen.

\subsection{Parametrisches Krümmungsmodell}
\begin{definition}[Parametrische Krümmung]
Ein Übergang wird beschrieben durch
\[
\kappa(s)=\kappa_0+(\kappa_1-\kappa_0)
\widehat\kappa\!\left(\frac{s}{L},\mathbf T\right),
\]
wobei \(L\) die Elementlänge, \(\mathbf T\) Übergangsparameter und
\(\widehat\kappa\) eine normierte Prototypfunktion bezeichnet.
\end{definition}

\subsection{Parametrisches Überhöhungsmodell}
\begin{definition}[Parametrische Überhöhung]
\[
\phi(s)=\phi_0+(\phi_1-\phi_0)
\widehat\phi\!\left(\frac{s}{L},\mathbf T_\phi\right).
\]
\end{definition}

\subsection{Parametervektor}
\begin{definition}[Reduzierter Parametervektor]
\[
\mathbf x=(L,\kappa_0,\kappa_1,\mathbf T,\phi_0,\phi_1,\mathbf T_\phi).
\]
\end{definition}
Die Dimension ist von der Diskretisierungsdichte unabhängig.

\subsection{Dimensionsreduktion}
Eine abgetastete Darstellung verwendet \(O(N)\) Variablen, die parametrische
Darstellung \(O(1)\) Parameter je Element. Dies reduziert die Problemgröße
drastisch.

\subsection{Ableitungsstruktur}
\begin{proposition}[Ableitungsstruktur]
\label{prop:parametric-derivative-structure}
Für einen Formparameter \(p\in\mathbf T\) bei festen \(L\), \(\kappa_0\)
und \(\kappa_1\) gilt
\[
\frac{\partial\kappa}{\partial p}
=(\kappa_1-\kappa_0)\frac{\partial\widehat\kappa}{\partial p}.
\]
Für die übrigen Parameter treten zusätzliche Terme auf:
\[
\frac{\partial\kappa}{\partial\kappa_0}=1-\widehat\kappa,
\qquad
\frac{\partial\kappa}{\partial\kappa_1}=\widehat\kappa,
\qquad
\frac{\partial\kappa}{\partial L}
=-(\kappa_1-\kappa_0)\,\frac{s}{L^2}\,
\partial_1\widehat\kappa\!\left(\frac{s}{L},\mathbf T\right),
\]
wobei \(\partial_1\) die Ableitung nach dem ersten Argument bezeichnet.
\end{proposition}
\begin{proof}
Gliedweise Differentiation des parametrischen Krümmungsmodells; die
\(L\)-Ableitung folgt aus der Kettenregel für \(s/L\).
\end{proof}
Alle Ableitungen werden analytisch aus Prototypfunktionen berechnet, sofern
der Parametertyp beachtet wird.

\subsection{Integration in die Rekonstruktion}
Die parametrische Darstellung wird bewertet durch
\[
\theta(s)=\int\kappa(s)\,\dd s,
\qquad\gamma(s)=\int(\cos\theta,\sin\theta)\,\dd s.
\]
So entsteht glatte, konsistente Geometrie ohne Diskretisierungsartefakte.

\subsection{Vergleich mit stichprobenbasierten Verfahren}
\begin{center}
\begin{tabular}{lcc}\toprule
& Stichprobenbasiert & Parametrisch\\\midrule
Dimension & hoch & niedrig\\
Glattheit & numerisch erzwungen & intrinsisch\\
Interpretierbarkeit & niedrig & hoch\\
Bedingungen & viele lokale & wenige globale\\\bottomrule
\end{tabular}
\end{center}

\subsection{Interpretation}
Der parametrische Ansatz verschiebt das Problem von der Anpassung von Werten
zum Entwurf von Funktionen.

\subsection{Verbindung zu Übergangsfamilien}
Verschiedene Übergangsbögen entsprechen verschiedenen Prototypfunktionen
\(\widehat\kappa\). Das Optimierungsproblem umfasst damit Auswahl und
Abstimmung von Übergangsfamilien.

\subsection{Rechnerische Folgen}
Es entstehen kleinere Probleme, bessere Konditionierung, schnellere Konvergenz
und unmittelbare Integration analytischer Jacobi-Matrizen.

\subsection{Verbindung zu klassischen Systemen}
Die Formulierung steht klassischen Systemen wie AXTRAN nahe, die auf
Elementparametern statt Stichprobendaten arbeiten.

\subsection{Verbindung zu AIM}
\begin{summary}
Die parametrische Darstellung ersetzt diskrete Abtastung durch strukturierte
Funktionen. Sie ermöglicht ein kompaktes, interpretierbares und analytisch
differenzierbares Modell als Grundlage effizienter, robuster
Alignment-Optimierung und ist zentral für AIM.
\end{summary}

\section{Hybrides parametrisch-diskretes Modell}
\subsection{Motivation}
Rein parametrische Modelle bieten niedrigdimensionale Struktur, können aber
für reale Daten zu unflexibel sein. Rein abgetastete Modelle sind flexibel,
führen jedoch zu großen, schlecht konditionierten Problemen. Ein hybrider
Ansatz verbindet beide Vorteile.

\subsection{Hybride Darstellung}
\begin{definition}[Hybrides Krümmungsmodell]
\[
\kappa(s)=\kappa_{\mathrm{param}}(s;\mathbf T)+\delta\kappa(s),
\]
wobei \(\kappa_{\mathrm{param}}\) das parametrische Basismodell und
\(\delta\kappa(s)\) einen Korrekturterm bezeichnet.
\end{definition}
\begin{definition}[Hybrides Überhöhungsmodell]
\[
\phi(s)=\phi_{\mathrm{param}}(s;\mathbf T_\phi)+\delta\phi(s).
\]
\end{definition}

\subsection{Diskretisierung der Korrektur}
\[
\delta\kappa_i=\delta\kappa(s_i),\qquad
\delta\phi_i=\delta\phi(s_i).
\]

\subsection{Erweiterter Parametervektor}
\begin{definition}
\[
\mathbf x=(\mathbf T,\mathbf T_\phi,
\delta\boldsymbol\kappa,\delta\boldsymbol\phi).
\]
\end{definition}

\subsection{Regularisierung der Korrekturen}
Zur Vermeidung von Überanpassung werden die Korrekturen bestraft:
\[
J_\delta=\int_0^L(\delta\kappa'(s))^2\,\dd s
+\int_0^L(\delta\phi'(s))^2\,\dd s.
\]
Dies erzwingt glatte, kleine Abweichungen vom parametrischen Modell.

\subsection{Residuenstruktur}
Das Residuum besteht aus parametrischen Beiträgen, Korrekturbeiträgen und
Datenanpassungstermen.

\subsection{Blockstruktur}
Der Parametervektor zerfällt in
\[
(\mathbf T,\delta\boldsymbol\kappa)
\quad\text{und}\quad
(\mathbf T_\phi,\delta\boldsymbol\phi),
\]
mit einer mehrstufigen Blockstruktur aus grobem Maßstab der Parameter und
feinem Maßstab der Korrekturen.

\subsection{Interpretation}
Das parametrische Modell erfasst geometrische Absicht, Übergangsstruktur und
Engineering-Semantik. Der Korrekturterm erfasst Messrauschen, lokale
Abweichungen und Modellierungsunvollkommenheiten.

\subsection{Verbindung zu Messdaten}
Messpolylinien können durch Minimierung von
\[
\sum_i\|\gamma(s_i)-p_i\|^2
\]
mit der hybriden Darstellung angepasst werden. Der parametrische Teil erklärt
die globale Struktur; die Korrektur nimmt Restfehler auf.

\subsection{Rechenstrategie}
Zunächst werden parametrische Variablen optimiert, danach Korrekturterme
verfeinert; optional können beide Ebenen abwechseln.

\subsection{Bezug zu LOD}
Das hybride Modell unterstützt natürlich Detaillierungsgrade: grober LOD nur
mit parametrischem Modell, feiner LOD mit Parametrik und Korrekturen.

\subsection{Verbindung zu AIM}
\begin{summary}
Die hybride Darstellung verbindet strukturierte Modellierung mit lokaler
Flexibilität. Sie ermöglicht robuste Alignment-Rekonstruktion aus
unvollkommenen Daten und bewahrt die Interpretierbarkeit parametrischer
Modelle. Damit verbindet sie Engineering-Entwurf und datengetriebene Anpassung
im AIM-Rahmen.
\end{summary}
